\section{Emit seven payload bits per byte}
Write the unsigned value in base 128. Each byte contributes seven low bits to the
number; its high bit indicates whether another byte follows. Emit the least significant
group first. For payloads \(p_0,\ldots,p_{\ell-1}\), reconstruct
\[
 u=\sum_{j=0}^{\ell-1}p_j128^j,\qquad0\leq p_j<128.
\]
All but the last byte have high bit one. The final byte has high bit zero. Zero needs
one byte, \texttt{00}; it is not represented by an empty sequence.

\begin{figure}[htbp]
\centering
\begin{tikzpicture}[alt={Unsigned 300 has base-128 payloads 44 and 2. The first byte has continuation bit one and payload binary 0101100, making hexadecimal AC. The second byte has continuation bit zero and payload binary 0000010, making hexadecimal 02. The value is 44 plus 128 times 2.},x=0.62cm,y=0.7cm,font=\small]
\foreach \row/\bits in {0/{1,0,1,0,1,1,0,0},2/{0,0,0,0,0,0,1,0}}{
 \foreach \bit [count=\j from 0] in \bits {
  \draw (\j,-\row) rectangle ++(1,-1);
  \node at (\j+0.5,-\row-0.5) {\bit};
 }
}
\draw[very thick,MARSAccent] (1,0)--(1,-1);
\draw[very thick,MARSAccent] (1,-2)--(1,-3);
\node at (0.5,0.6) {More?};
\node at (4.5,0.6) {Seven payload bits};
\node[anchor=west] at (8.4,-0.5) {\texttt{ac}: payload 44};
\node[anchor=west] at (8.4,-2.5) {\texttt{02}: payload 2};
\node at (4,-3.8) {\(44+128\cdot2=300\)};
\end{tikzpicture}
\caption{The two bytes of unsigned 300. Byte order follows payload significance.}
\end{figure}

This grouping is not the byte order of a fixed-width memory word. For example,
little-endian storage of the 32-bit integer 300 is \texttt{2c 01 00 00}; this
variable-length representation is \texttt{ac 02}. The high bit participates in framing,
so it must be masked off before shifting a payload into the accumulator.

The unsigned grouping is commonly described as unsigned LEB128 or a base-128 varint.
This project applies it after the signed mapping above. Signed LEB128 is another
representation; its sign-extension rules do not describe these mapped payloads.
The complete project format is defined by this chapter, including its strict
minimal-length rule and external input extent.
